% Pursuing Stacks (cotes 134-2 to 134-8): archive cote 134-7, site page 16 (PDF page 17), typescript page 458.
% typescript. Draft compared: ps6.tex:735-796.
% AI-transcribed by Claude Opus 5.5 (claude-opus-5-5, adaptive thinking, effort high), Message Batches API, batch ps3 (msgbatch_01XdFmuVVBAcgc7EFxdtCqsj), 2026-11-12; not yet checked by eye.
\scanpage{134-7/16}
\mspage{458}

varying $M_{*}$ , is that to a direct sum corresponds the componentwise
tensor product \uline{over} $\underline{Z}$ , whereas for the last term of(w) we have to take
tensor products \uline{over} k (in the middle term, \uline{completed} tensor products
over k) - the two variances are clearly at odds which each other.

Thus, when working with would-be ``k-homotopy types'' as defined
here via semisimp\unsure{l|h}icial objects in $U(k)^{\bullet}$ , \uline{we should altogether drop the}
\uline{idea that the homology groups of the corresponding semisimplicial set}
$X_{*}(k)$ \uline{are k-modules}. I would'nt really look at them as being ``the''
homology groups of the k-homotopy type $X_{*}$ , these should be rather given
via k-linearization \xout{P} $L_{pt}$ and they are k-modules, indeed they come
from \over{an}{a} \inserted{canonical} object of $D.(Ab_{k})$ , namely $L_{pt}(X_{*})$ . In the example just
looked at, presumably we should get the exterior algebra of M \uline{over} k
(not $\underline{Z}$ !) . This makes me suspect even that, except in the case $k=\underline{Z}$ ,
this semisimplicial set $X_{*}(k)$ does'nt make much sense, namely it\unsure{\deleted{'}}s
homology invariants (and presumably it's homotopy groups too) are not
really relevant for the k-homotopy type $X_{*}$ , \uline{which has invariants of it's}
\uline{own which are completely different}. \uline{Thus I am not} \xout{all} \uline{at all convinced}
\margin{X} \uline{any more that the homotopy groups} $\hand{\pi}_{i}(X_{*})$ \uline{carry k-module structures} (as
expected at the beginning, p. 446) - but to clear our mind on that
matter, we should take off from the simplistic example when $X_{*}$ comes
from a semisimplicial k-module $M_{*}$ , in which case of course we have by
Whitehead's isomorphism
\[ \hand{\pi}_{i}(X_{*}(k)) \hand{\simeq} H_{i}(M_{\unsure{*}}.) \quad ( \hand{\simeq} \hand{\pi}_{i}(M_{*})) \quad , \]
where M. is the chain complex associated to $M_{*}$ . Also, \deleted{I am not com-}\inserted{it is now}
\deleted{pletely clear still} \inserted{becoming dubious} whether weak equivalences \xout{\unsure{in}} for objects $X_{*}$ should
be defined (as we did) via $X_{*}(k)$, as suggested by the \xout{example} ``simplis-%
tic'' example above. \inserted{If we} \over{D}{d}efin\unsure{\over{ing}{e}} homology of $X_{*}$ by the formula
\[ \text{(z)} \qquad LH.(X_{*}) \overset{\text{def}}{=} L_{pt}(X_{*}) \text{ viewed as an object of } D.(Ab_{k}) \ , \]
and, accordingly the homology modules
\[ \text{(z')} \qquad H_{i}(X_{*}) \overset{\text{def}}{=} \hand{\pi}_{i}(L_{pt}(X_{*})) = H_{i}(LH.(X_{*})) \quad , \]
maybe the better idea for defining weak equivalences $X_{*} \hand{\longrightarrow} X'_{*}$ is
by demanding that \uline{they should be transformed into quasi-isomorphisms by}
\uline{the total homology functor} LH. , or equivalently, induce isomorphisms
for the homology modules (z'). If the answer to our crucial question
is affirmative (with the corrections made, including $k = \underline{Z}$ ), then

\margin{16}
